\documentclass[11pt,a4paper]{article}
\usepackage{fontspec,amsmath,amssymb,amsthm,microtype}
\setmainfont{lmroman10-regular.otf}[BoldFont=lmroman10-bold.otf,ItalicFont=lmroman10-italic.otf,BoldItalicFont=lmroman10-bolditalic.otf]
\usepackage[margin=25mm]{geometry}
\usepackage[unicode,hidelinks]{hyperref}
\hypersetup{pdfauthor={LangevinLean Project}}
\setlength{\emergencystretch}{3em}
\setlength{\parskip}{0.4em}
\setlength{\parindent}{0pt}
\allowdisplaybreaks[2]
\raggedbottom
\title{\texorpdfstring{Equality of initial position and common-\allowbreak{}velocity phase transport distances}{Equality of initial position and common-velocity phase transport distances}}
\author{LangevinLean Project}
\date{}
\begin{document}\maketitle
Draft mathematical exposition; these notes may contain mistakes.


The common initial velocity distribution is independent of each initial position distribution. Adjoining that identical velocity does not change their exact quadratic transport distance. This identifies the position-distance estimate with the initial phase-product distance printed in the source theorem, for the Brownian-driven solution.



Appending the same independent velocity gives an upper transport bound, while position projection gives the reverse bound. Together they identify the initial phase-product distance exactly with position Wasserstein distance.


\section{\texorpdfstring{The exact product transport identity}{The exact product transport identity}}


For the upper bound, couple the two initial positions and use the identical independent velocity in both phase states. Their phase displacement has zero velocity component, so its norm is exactly the original position displacement. The random-map transport theorem then gives the upper bound. For the reverse inequality, the position projection has Lipschitz constant one. Projecting every phase coupling therefore gives the position transport bound. The two extended-real inequalities prove exact equality, even before imposing finite moments. Real conversion preserves this equality.



For the lower bound, project any phase coupling onto positions. For the upper bound, extend any position coupling by one identical independent velocity. The added coordinate contributes zero cost. Taking both infima yields equality of the phase-product and position transport distances.



Earlier results used here: \cite[Section~3.17.3]{LangevinLean2026ULASC}; \cite[Section~3.17.4]{LangevinLean2026ULASC}.

\section{\texorpdfstring{The source theorem’s literal full phase-\allowbreak{}product right-\allowbreak{}hand side}{The source theorem’s literal full phase-product right-hand side}}


Apply the already proved specified Brownian mixing estimate, with arbitrary common finite-second-moment velocity distribution. The source restriction \(0\le a<\gamma/2\) gives the required positive denominator. Replacing the initial position Wasserstein distance by the exact common-velocity phase-product identity produces the source position-contraction formula with precisely its printed prefactor, exponent and initial phase-product distance. The Gibbs corollary specializes to Gaussian velocity and the proved stationary target, recovering the explicit \(\sqrt2 e^{-mt/\gamma}\) friction-regime statement with the literal initial Gaussian-times-position phase distance.



The theorem’s initial phase distribution has the same Gaussian velocity factor as the stationary target. The product transport identity therefore rewrites its phase distance exactly as the initial position distance. No extra velocity discrepancy is hidden in the printed right-hand side.



Earlier results used here: \cite[Section~3.17.4]{LangevinLean2026ULASC}; \cite[Section~3.26.13]{LangevinLean2026ULASC}; \cite[Section~3.26.3]{LangevinLean2026ULASC}.

\begin{thebibliography}{ULA-LSI}
\bibitem[ULA-SC]{LangevinLean2026ULASC}
LangevinLean Project. \emph{Unadjusted Langevin under Strong Convexity}. LL-01, 2026. Draft report, content revision 7e430ee07530. \url{https://odakyildiz.github.io/LangevinLean/docs/pdf/UnadjustedLangevin/FullProof.pdf}.
\end{thebibliography}

\end{document}
